% R1 pages of the R1--R3 exhibit book (input into r123_exhibit_book.tex).
% Every number is taken from the block's tables/*.csv, INDEX.md or a review's re-check.
\section{R1 --- How internationalised firms responded to the Paris Agreement}

{\footnotesize\color{soft}Block folder: \texttt{20260917-run/block-r1-paris-internationalisation/} (PIC Toar). Deasy's row, verbatim:}
\begin{quote}\footnotesize
\textbf{Request.} ``How internationalized firms respond to climate commitment?? Do they adjust their energi use, emisi, dan juga
internationalization strategy-nya? Do global climate commitments reach the factory floor? Did internationally engaged firms
change their environmental behaviour more strongly after the Paris Agreement? Outcome variables: VA/energy use, scope 1, scope 2,
total; VA/emisi, scope 1, scope 2, total; Energi use, scope 1, 2, total; Emisi, scope 1, 2, total; GJ/emisi, scope 1, 2, total;
Co2/va; Gj/va; Electricity; Self generation.''\par
\textbf{Strategy.} ``Matching per year: export vs non export, import vs non import, fdi vs non fdi. DID. Apakah mungkin dengan
staggered did?? Alternative: Freeze internationalisation status sebelum Paris, misalnya average 2013--14: pre exporters x year;
2015 transition year; 2016+ post; Firm FE + industry$\times$year FE + province$\times$year FE. Matching/entropy balancing dapat
digunakan sebagai robustness untuk membuat pre-Paris groups comparable.''\par
\textbf{Notes.} ``Sesuai diskusi jumat (28 august), karena kita gak bs tau first time exporting/importing/fdi nya, kita bs coba
dengan method matching yearly.. dan dicompare bagaimana behaviour mereka toward energy use, emisi, ataupun ekspor, impor, fdi\ldots''
\end{quote}

{\footnotesize
\renewcommand{\arraystretch}{1.08}
\begin{tabularx}{\linewidth}{@{}l X p{4.3cm} p{4.6cm}@{}}
\toprule
ID & Deasy's item (short) & Exhibits & Status \\
\midrule
Q1 & How internationalized firms respond to climate commitment & R1-02 to R1-10 & done on this data \\
Q2 & Do they adjust their energi use & R1-05, R1-06, R1-07, R1-09 & done \\
Q3 & \ldots emisi & R1-05, R1-06, R1-09 & done \\
Q4 & \ldots internationalization strategy-nya & R1-01; R2-01 (start years) & done, descriptive \\
Q5 & Do global climate commitments reach the factory floor? & R1-05, R1-06, R1-09, R1-10 & partial: Paris came to all firms at once \\
Q6 & Did internationally engaged firms change more strongly after Paris? & R1-02 to R1-10 & done on this data \\
O1a--c & VA/energy use: scope 1, 2, total & R1-05, R1-09; total also R1-02 to R1-04, R1-08 & done (reciprocals, in tables) \\
O2a--c & VA/emisi: scope 1, 2, total & R1-05, R1-09; total also R1-02 to R1-04, R1-08 & done (reciprocals, in tables) \\
O3a--c & Energi use: scope 1, 2, total & R1-05, R1-09; scope 2 also R1-07 & done \\
O4a--c & Emisi: scope 1, 2, total & R1-05, R1-09 & done; scope 2 repeats scope-2 energy \\
O5a--c & GJ/emisi: scope 1, 2, total & R1-05, R1-09 & done; scope 2 is the same for every firm \\
O6, O7 & Co2/va; Gj/va & R1-02 to R1-06, R1-08, R1-09 & done \\
O8 & Electricity & R1-05, R1-07, R1-09 & done \\
O9 & Self generation & R1-05, R1-07, R1-09 & done; 2014--15 too few firms recorded \\
G1 & Export vs non export & R1-01, R1-02, R1-05 to R1-10 & done; 2022 has no known non-exporters \\
G2 & Import vs non import & R1-01, R1-03, R1-05 to R1-10 & done; import share jumps in 2017 \\
G3 & Fdi vs non fdi & R1-01, R1-04, R1-05 to R1-10 & done; 2020 spike left out \\
S1 & Matching per year & R1-08 & done \\
S2 & DID & R1-09 & done, association only \\
S3 & Apakah mungkin dengan staggered did?? & R2 block (R2-01, R2-03 to R2-09) & done in R2; Paris is one date for all \\
S4 & Freeze status before Paris (2013--14): pre exporters x year & R1-06, R1-09, R1-10 & done \\
S5 & 2015 transition year & R1-06, R1-09 & done \\
S6 & 2016+ post & R1-06, R1-09 & done: 2018 drawn, 2017--22 in table \\
S7 & Firm FE + industry$\times$year FE + province$\times$year FE & R1-09 & done \\
S8 & Matching/entropy balancing as robustness & R1-10 & done \\
N1 & Matching yearly; compare energy use, emisi, ekspor, impor, fdi & R1-01 to R1-04, R1-07, R1-08 & done \\
\bottomrule
\end{tabularx}\par}
{\scriptsize\color{soft}IDs are \texttt{R1.}\emph{ID} in TRACKER.md. Scope 1 = fuels burned; scope 2 = bought electricity. DID = the group's change
minus other firms' change. FE = each firm, industry-year and province-year measured against its own average. Staggered = firms
starting in different years.\par}

\medskip
{\small\textbf{What this block says on the development data.}
\begin{itemize}
\item Exporters, importers and foreign-owned firms mostly use less energy and CO2 per unit of value added than other firms,
  before and after Paris (R1-02 to R1-04). Set against firms of the same industry and size, only importers stay clearly lower in most years (R1-08).
\item From 2013--14 to 2018 no group changed its energy or CO2 per unit of value added clearly differently from other firms:
  not with groups fixed before Paris (R1-06), not matched (R1-08), not in the regression (R1-09), not reweighted (R1-10).
\item Exporters used 27\% less energy than similar firms by 2018 (R1-09), but their value added fell 25\%; reweighted, the gap is
  $-17$\% and not clear (R1-10). That points to size, not to a response to Paris.
\item Several gaps were already opening before Paris (R1-09: 18 of 51 before-Paris comparisons clear, 10 of 51 after).
  Nothing here links a change to the Paris Agreement.
\item Recording problems shape the later years: far more importers from 2017, a one-year foreign-ownership spike in 2020,
  energy records that jump in 2017 and drop in 2019--22. 2018 is the one solid year after Paris.
\item Firm links are fictional, so the results that follow firms (R1-06, R1-09, R1-10) test the method; re-read them on the real files.
\end{itemize}}

% ---------------------------------------------------------------------------------------------------------------
\expage{R1-01}{Who exports, imports or is foreign-owned, by year}{\blockI/figures/01_group_shares_support.pdf}{\texttt{01\_group\_shares\_support.do} $\to$ \texttt{figures/01\_\ldots.png}, \texttt{tables/01\_group\_shares\_support.csv}, \texttt{tables/01\_group\_support\_cells.csv}}{0.60}
\expart{What it is}{For each year 2010--2022, the share of firms that export (a), import (b) or have 10\% or more foreign
ownership (c), among firms whose status is known. The rows under each panel count the firms in the group, not in it, and of
unknown status. These counts are the group sizes behind every later R1 comparison.}
\expart{How to read it}{Each dot is one year's share; the first and last shares are printed. The gray band at 2016 marks the
missing annual file, so the line breaks there. Firms of unknown status (gray count row) are left out of the share. No 2022 export
share is drawn: the 2022 file records only direct exports, so no firm is known to be a non-exporter.}
\expart{What it shows on this data}{Exporters are 15--30\% of firms with a known export status and foreign-owned firms 8--11\%,
with no lasting change after Paris. Importers rise from 17\% (2010) to 26\% (2015), then jump to 56\% in 2017; foreign ownership spikes
to 47\% in 2020 only. Both jumps more likely reflect how the survey recorded these items than what firms did. Export status is
unknown for 29--57\% of firms each year.}
\expart{How it helps Deasy's request}{It answers ``\ldots dan juga internationalization strategy-nya?'' descriptively, year by
year; each firm's start year is in R2-01. It sets up the three comparisons of the strategy line, ``export vs non export, import vs
non import, fdi vs non fdi'', and shows how many firms each has. It cannot show which firms changed strategy because of Paris.}
\items{R1.Q4, R1.G1, R1.G2, R1.G3, R1.N1}
\expart{Keep in mind}{The firms surveyed change from year to year and the chart follows no firm. Foreign ownership is drawn at 10\%
or more; any foreign ownership and the 2022 direct-export question (344 yes, 1,988 no, 845 unknown) are in the table only.}

% ---------------------------------------------------------------------------------------------------------------
\expage{R1-02}{Exporters against non-exporters, year by year}{\blockI/figures/02_exporters_energy_co2_per_va.pdf}{\texttt{02\_exporters\_energy\_co2\_per\_va.do} $\to$ \texttt{figures/02\_\ldots.png}, \texttt{tables/02\_\ldots.csv}}{0.58}
\expart{What it is}{The middle (median) firm's energy (GJ of fuels and bought electricity) and CO2 (tonnes) per million of value
added, exporters against non-exporters, each year 2010--2021. The bottom panels divide the exporter line by the non-exporter line,
with a 95\% range from 999 random redraws of each year's firms.}
\expart{How to read it}{Top: black dots = exporters, gray squares = non-exporters; first and last values are printed. Bottom: a
diamond below 100 means exporters need less energy or CO2 per unit of value added (not less in total: they are bigger firms); where
its vertical line crosses 100 there is no clear difference. The dashed line marks the Paris Agreement (December 2015); the gray
band is 2016, which has no annual file. 2022 has no known non-exporters, so no comparison is drawn.}
\expart{What it shows on this data}{Exporters' middle firm uses 54--101\% of non-exporters' energy per unit of value added. The gap
is clear in 2010 (68) and 2014 (67) before Paris, and in 2018 (73), 2019 (54) and 2021 (67) after; 2017 sits at 100. CO2 follows
the same pattern but is clear less often (2018: 81, range 64--107) and in 2017 is clearly above 100 (129). After Paris the gap is clear in more years, but it was already
there before and jumps from year to year.}
\expart{How it helps Deasy's request}{It answers ``Did internationally engaged firms change their environmental behaviour more
strongly after the Paris Agreement?'' for exporters, with ``Gj/va'' and ``Co2/va'' (total) drawn. ``VA/energy use'' and ``VA/emisi''
(total) are the same medians turned upside down, in the table. It is the unmatched yearly comparison; the matched one asked for in
``Matching per year'' is R1-08.}
\items{R1.Q1, R1.Q6, R1.O6, R1.O7, R1.O1c, R1.O2c, R1.G1, R1.N1}
\expart{Keep in mind}{Different firms each year (none is followed). Exporters are about four times larger (median about 150 workers
against 40) and in other industries; firms of unknown export status (29--58\% a year) are left out. Value added is at current
prices, so price rises alone lower every line. Association, not cause.}

% ---------------------------------------------------------------------------------------------------------------
\expage{R1-03}{Importers against non-importers, year by year}{\blockI/figures/03_importers_energy_co2_per_va.pdf}{\texttt{03\_importers\_energy\_co2\_per\_va.do} $\to$ \texttt{figures/03\_\ldots.png}, \texttt{tables/03\_\ldots.csv}}{0.60}
\expart{What it is}{The same chart as R1-02 for importers against non-importers: the middle firm's energy (GJ) and CO2 (tonnes)
per million of value added each year, and importers as \% of non-importers with a 95\% range from 999 redraws.}
\expart{How to read it}{As R1-02: below 100 = importers need less energy or CO2 per unit of value added; a vertical line crossing
100 = no clear difference; gray band = no 2016 file. 2017 is not drawn: importers' recorded energy jumps that year (median 1,404
against 157 GJ per million of value added). Hollow gray marks (2019--22) are years whose importer energy records are in doubt.}
\expart{What it shows on this data}{Before Paris, importers' middle firm uses 60--87\% of non-importers' energy per unit of value
added, clearly less in every year except 2012. In 2018, the one later year that can be read, it is 59\% (range 53--65), about the
2010--11 gaps (62 and 60); CO2 is 56\%. From 2019 the gap drops to 14--25\%, but about 1,300 importers with no export answer
report a fifth of their earlier fuel and electricity per worker while their value added per worker stays flat: a recording
problem, not behaviour.}
\expart{How it helps Deasy's request}{It covers ``import vs non import'' and the after-Paris question for importers, with ``Gj/va''
and ``Co2/va'' drawn and ``VA/energy use'', ``VA/emisi'' (total) in the table. It answers only in part: 2018 is the one readable
later year, and by then the files count far more firms as importers (26\% in 2015, 56\% in 2017).}
\items{R1.Q1, R1.Q6, R1.O6, R1.O7, R1.O1c, R1.O2c, R1.G2, R1.N1}
\expart{Keep in mind}{From 2017 the importer group is a different mix of firms (likely a recording change), so a later gap may
reflect who is counted. Firms differ each year and the groups differ in size and industry. Association, not cause.}

% ---------------------------------------------------------------------------------------------------------------
\expage{R1-04}{Foreign-owned against other firms, year by year}{\blockI/figures/04_foreign_owned_energy_co2_per_va.pdf}{\texttt{04\_foreign\_owned\_energy\_co2\_per\_va.do} $\to$ \texttt{figures/04\_\ldots.png}, \texttt{tables/04\_\ldots.csv}}{0.60}
\expart{What it is}{The same chart for firms with 10\% or more foreign ownership against other firms: the middle firm's energy (GJ)
and CO2 (tonnes) per million of value added each year, and foreign-owned as \% of other firms with a 95\% range from 999 redraws.}
\expart{How to read it}{As R1-02: below 100 = less energy or CO2 per unit of value added; crossing 100 = no clear difference; gray
band = no 2016 file. 2020 is not drawn: 47\% of firms are marked foreign-owned that year, against 8--11\% otherwise. Light-gray
marks in 2017, 2019, 2021 and 2022, in both groups, are years when many firms' energy records jump or drop.}
\expart{What it shows on this data}{In 2010--14 the foreign-owned middle firm used 47--76\% of other firms' energy per unit of
value added, clearly less in every year. By 2015 the gap had closed (101). In 2018 it is 79\% (range 62--95), clear; CO2 is 91\%
(76--105), not clear. The gap narrowed before Paris, not after it.}
\expart{How it helps Deasy's request}{It covers ``fdi vs non fdi'' and the after-Paris question for foreign-owned firms, with
``Gj/va'' and ``Co2/va'' drawn and ``VA/energy use'', ``VA/emisi'' (total) in the table. It answers in part: 2018 is the one clean
later year.}
\items{R1.Q1, R1.Q6, R1.O6, R1.O7, R1.O1c, R1.O2c, R1.G3, R1.N1}
\expart{Keep in mind}{Foreign-owned firms are much larger (median 141--201 workers against 38--53). Unknown ownership (up to 10\%
of firms, in 2021) is left out. Firms differ each year. Association, not cause.}

% ---------------------------------------------------------------------------------------------------------------
\expage{R1-05}{Before against after Paris, every outcome}{\blockI/figures/05_before_after_paris.pdf}{\texttt{05\_before\_after\_paris.do} $\to$ \texttt{figures/05\_\ldots.png}, \texttt{tables/05\_before\_after\_paris.csv}}{0.57}
\expart{What it is}{For every outcome in Deasy's list, the middle (median) firm before Paris (average of 2013 and 2014) against
after (2018): energy and CO2 per unit of value added, energy used, CO2 emitted and GJ per tonne of CO2, each for fuels burned
(scope 1), bought electricity (scope 2) and the total; bought and own electricity; value added and workers for comparison.
Panel (a) is all firms; (b)--(d) set exporters, importers and foreign-owned firms against the other firms.}
\expart{How to read it}{In (a) the dot is the \% change, its line the 95\% range, and 0 = no change. In (b)--(d) the black dot is
the group's change and the square the other firms'; the number is the gap (group's change relative to the other firms', 0 = both
changed alike). Black links and numbers are clear; gray ones are not.}
\expart{What it shows on this data}{All firms use less energy per unit of value added ($-29$\%; fuels $-46$\%) but more energy in
total (+66\%), as value added rose 155\%. Much of this is who is counted: importers with no export answer grow from 12--14\% of firms
in 2013--14 to 41\% in 2018; without them energy per value added falls 13\% and value added rises 62\%. Exporters show no clear gap
except more firms making their own electricity (+40). Importers' clear gaps (energy per value added $-21$, CO2 $-38$) sit mostly in
bought electricity. Foreign-owned firms' lower energy use ($-44$) goes with much slower value-added growth ($-62$).}
\expart{How it helps Deasy's request}{It answers ``Do global climate commitments reach the factory floor?'' in part: it shows what
changed from 2013--14 to 2018, not whether Paris caused it. It covers the whole outcome list. GJ per tonne of bought-electricity CO2
is the same for every firm ($-2$\%, one grid factor a year); ``VA/energy use'' and ``VA/emisi'' are in the table.}
\items{R1.Q2, R1.Q3, R1.Q5, R1.Q6, R1.O1a--O9, R1.G1, R1.G2, R1.G3}
\expart{Keep in mind}{After = 2018 only: other later files have energy records that jump or drop. Most rows were already moving
this way before Paris (energy per value added $-34$\% from 2010 to 2013--14). Own-electricity rows use 2013 as ``before''. Firms
differ each year.}

% ---------------------------------------------------------------------------------------------------------------
\expage{R1-06}{Groups fixed before Paris, followed over time}{\blockI/figures/06_frozen_pre_paris_groups.pdf}{\texttt{06\_frozen\_pre\_paris\_groups.do} $\to$ \texttt{figures/06\_\ldots.png}, \texttt{tables/06\_frozen\_pre\_paris\_groups.csv}}{0.58}
\expart{What it is}{Each firm is put once into a group from its 2013--14 records (exporter, importer or 10\%+ foreign-owned in either
year), then followed 2010--2022. Lines show the middle firm's energy and CO2 per unit of value added as \% of the group's own
2013--14 level. Groups: 592 exporters against 1,727 other firms, 891 importers against 2,140, 397 foreign-owned against 2,634.
Left out: 7,432 firms with no 2013--14 record and 712 with no known export answer.}
\expart{How to read it}{Both lines average 100 in 2013--14, so the chart compares changes, not levels. The top-right text is the
2018 difference, group line minus other firms' line in points, with a 95\% range from 999 redraws of whole firms (gray = includes
0). 2015 is the transition year; 2017 is not drawn (energy records jump); 2019--22 are light gray (records in doubt); the gray band
is 2016, which has no file.}
\expart{What it shows on this data}{By 2018 each group fell to about 65--80\% of its 2013--14 level, as other firms did. The 2018
differences are $-9$ and $-4$ points for exporters (energy, CO2), $-2$ and 0 for importers, $-2$ and $-9$ for foreign-owned firms;
none is clear. The paths were not parallel before Paris: in 2010--12 importers stood 39--47 points above other firms (falling
faster, clear). Counting only firms in their group in every known 2013--14 year, foreign-owned firms fall 32 points more (range
$-48$ to $-9$), clear, but on 43 firms.}
\expart{How it helps Deasy's request}{It draws Deasy's ``Freeze internationalisation status sebelum Paris, misalnya average 2013--14:
pre exporters x year; 2015 transition year; 2016+ post''; the regression version is R1-09. ``2016+'' is read at 2018, the one later
year whose records look normal. For the after-Paris question: no group moved clearly differently by 2018.}
\items{R1.S4, R1.S5, R1.S6, R1.Q2, R1.Q3, R1.Q5, R1.Q6, R1.O6, R1.O7, R1.G1, R1.G2, R1.G3}
\expart{Keep in mind}{Firm links are fictional: in 2018 only 39\% of answering exporters, 71\% of importers and 15\% of
foreign-owned firms are still so recorded. The answer depends on the grouping rule. Change, not its cause.}

% ---------------------------------------------------------------------------------------------------------------
\expage{R1-07}{Bought and own electricity, by year and group}{\blockI/figures/07_electricity_own_generation.pdf}{\texttt{07\_electricity\_own\_generation.do} $\to$ \texttt{figures/07\_\ldots.png}, \texttt{tables/07\_electricity\_own\_generation.csv}}{0.60}
\expart{What it is}{Year by year, exporters, importers and 10\%+ foreign-owned firms, each against other firms of known status:
(a)--(c) the middle firm's electricity bought from PLN (the state power company) and other suppliers, in thousand kWh per million of value added; (d)--(f)
the share of firms making some of their own electricity; (g)--(i) own electricity as \% of the electricity used, for the middle
firm among those making some.}
\expart{How to read it}{Black dots = the group, squares = other firms; first and last values are printed. Hollow gray marks are
years whose records are in doubt (as in R1-03 and R1-04). ``Not drawn'' under each panel lists years left out: 2014--15 for own
electricity (recorded for only 56--60\% of firms), 2022 for exporters, 2020 for foreign-owned firms. Gaps and 95\% ranges are in
the table.}
\expart{What it shows on this data}{Exporters bought 118--175\% of non-exporters' electricity per unit of value added before Paris,
112\% (range 80--151) in 2018 and 93--104\% in 2019--21: the gap closes, but not clearly. Exporters make their own electricity more
often: +7 to +14 points before Paris, +23 in 2018. Importers' 2018 lead in making their own (78\% against 16\%) is mostly who the
files count as importers: in 2017, 90\% of importers have no export answer, and 77\% of firms with no export answer make some of
their own.}
\expart{How it helps Deasy's request}{It answers ``Electricity'', ``Self generation'' and ``Energi use \ldots\ scope 2'' (bought
electricity) year by year for the three groups; amounts per firm are in the table. With R1-05 and R1-09 it covers these items. It
shows gaps, not why they changed.}
\items{R1.O8, R1.O9, R1.O3b, R1.Q2, R1.Q6, R1.G1, R1.G2, R1.G3, R1.N1}
\expart{Keep in mind}{From 2017 makers record a far smaller own share in every group (mostly 1--8\%), likely a recording change,
so compare groups within a year. Firms differ each year; value added is at current prices.}

% ---------------------------------------------------------------------------------------------------------------
\expage{R1-08}{Each group against similar firms, year by year (matching)}{\blockI/figures/08_matched_gaps_by_year.pdf}{\texttt{08\_matched\_gaps\_by\_year.do} $\to$ \texttt{figures/08\_\ldots.png}, \texttt{tables/08\_matched\_gaps\_by\_year.csv}}{0.57}
\expart{What it is}{Each year, each exporter (importer, foreign-owned firm) is compared only with other firms of the same year,
2-digit industry and size band (under 50, 50--99, 100--249, 250--499, 500+ workers); 92--100\% of group firms find such firms.
Diamond = the typical group firm's energy or CO2 per unit of value added as \% of its similar firms', with a 95\% range from 999
redraws of whole firms. Circle = the same against all other firms (not matched).}
\expart{How to read it}{100 = same as the firms compared; a range crossing 100 = no clear difference. Diamond against circle shows
what matching changes. ``Typical'' is an average in which a firm at half and one at double cancel out, so circles differ from R1-02
to R1-04's medians. Top right: 2013--14 and 2018 values and the change in points (gray = not clear). Importers 2017 and
foreign-owned 2020 are not drawn; light gray = records in doubt.}
\expart{What it shows on this data}{Matched, exporters use about as much as similar non-exporters (70--125 in 2010--15). Importers
use less (70--90 before Paris, clear in every year but 2012; 70 energy and 67 CO2 in 2018, clear). Foreign-owned firms are at 63--83 in 2010--14 and about level
from 2015. From 2013--14 to 2018 no change is clear: exporters +1 and 0 points, importers $-7$ and $-12$, foreign-owned +17 and +19.
Without matching, importers' falls would be clear ($-15$ and $-20$).}
\expart{How it helps Deasy's request}{This is ``Matching per year: export vs non export, import vs non import, fdi vs non fdi'' and
the notes' ``kita bs coba dengan method matching yearly.. dan dicompare bagaimana behaviour mereka toward energy use, emisi'', done
for ``Gj/va'' and ``Co2/va'' (``VA/energy use'', ``VA/emisi'' in the table). Matched, no group changed clearly after Paris.}
\items{R1.S1, R1.N1, R1.Q1, R1.Q6, R1.O6, R1.O7, R1.O1c, R1.O2c, R1.G1, R1.G2, R1.G3}
\expart{Keep in mind}{Matching evens out only industry and size, and it does not fix the doubtful records. Firms differ each year;
differences, not their cause.}

% ---------------------------------------------------------------------------------------------------------------
\expage{R1-09}{Groups fixed before Paris, as a regression}{\blockI/figures/09_frozen_groups_regression.pdf}{\texttt{09\_frozen\_groups\_regression.do} $\to$ \texttt{figures/09\_\ldots.png}, \texttt{tables/09\_frozen\_groups\_regression.csv}}{0.57}
\expart{What it is}{R1-06's groups (fixed once from 2013--14 records), for every outcome in Deasy's list. The black dot is how much
more the group's measure changed from 2013--14 to 2018 than that of firms in the same industry and province, each firm measured
against its own average: a \% difference (electricity rows in points) with a 95\% interval.}
\expart{How to read it}{0 = the group changed like similar firms; black number = clear, gray = not. The gray triangle is the same
for 2010--12 against 2013--14: away from 0 means the group was already moving apart before Paris. The hollow circle compares with
all other firms, without the same-industry and same-province matching. Read energy used together with the value added and workers
rows: a shrinking group uses less energy.}
\expart{What it shows on this data}{Energy and CO2 per unit of value added: no group differs clearly (total energy: exporters
$-3$\%, importers +8\%, foreign-owned +10\%). Exporters used 27\% less energy (interval $-43$ to $-7$) and emitted 29\% less CO2
than similar firms, but their value added fell 25\% and workers 14\%: this points to size, not to a Paris response. 10 of 51 after-Paris
markers are clear, against 18 of 51 before: in 2010--12 exporters' and importers' energy per value added were already +31\% and
+27\%.}
\expart{How it helps Deasy's request}{It builds ``DID'' (the group's change minus similar firms' change) on the frozen groups with
``Firm FE + industry$\times$year FE + province$\times$year FE''. The ``2015 transition year'' and the literal ``2016+ post'' (2017--22
against 2010--14) are in the table. In 2015 importers' energy per value added is +28\%, clear; for ``2016+'' no total
differs clearly. Every listed outcome is covered.}
\items{R1.S2, R1.S4, R1.S5, R1.S6, R1.S7, R1.Q1, R1.Q2, R1.Q3, R1.Q5, R1.Q6, R1.O1a--O9, R1.G1, R1.G2, R1.G3}
\expart{Keep in mind}{Association, not the effect of Paris. Counting only firms in their group in every known 2013--14 year, 16 of
51 are clear (foreign-owned energy used $-54$\%, value added $-36$\%): the rule matters. Firm links are fictional (39\%, 71\%, 15\%
still so recorded in 2018).}

% ---------------------------------------------------------------------------------------------------------------
\expage{R1-10}{Groups fixed before Paris, other firms reweighted}{\blockI/figures/10_reweighted_frozen_groups.pdf}{\texttt{10\_reweighted\_frozen\_groups.do} $\to$ \texttt{figures/10\_\ldots.png}, \texttt{tables/10\_reweighted\_frozen\_groups.csv}, \texttt{tables/10\_\ldots\_balance.csv}}{0.60}
\expart{What it is}{R1-09 again, with the other firms reweighted (entropy balancing): each gets one weight so that in 2013--14 the
weighted other firms match the group in workers, value added, energy and CO2 per unit of value added, bought-electricity share, own
electricity, industry and island. Top: the 2013--14 gaps. Bottom: the change from 2013--14 to 2018, group minus reweighted other
firms, with a 95\% interval.}
\expart{How to read it}{Top, in spreads between firms (standard deviations): $\circ$ = before weighting, $\bullet$ = after (0 =
matched), $\diamond$ = among firms also recorded in 2018, on which the bottom rows rest. ``Worth'' = how many equally counted firms
the weighted other firms amount to (283, 338, 570). Bottom, as R1-09: $\bullet$ reweighted, $\circ$ without weights, gray triangle
= 2010--12 check; gray number = not clear.}
\expart{What it shows on this data}{Before weighting the groups were 0.92--1.13 spreads larger in workers and value added; among
firms still recorded in 2018 exporters are again larger (workers +0.39). Reweighted, 2 of 51 after-Paris markers are clear (10
without). Exporters' energy used is $-17$\% ($-36$ to +7), no longer clear. Importers' value added grows 34\% more (+6 to +71,
clear), so their energy per value added falls ($-25$\%, not clear) while energy used does not (+1\%). Before-Paris gaps disappear
(2 of 51 clear, 18 without).}
\expart{How it helps Deasy's request}{It is ``Matching/entropy balancing dapat digunakan sebagai robustness untuk membuat pre-Paris
groups comparable'', applied to R1-09's design and built in Stata with no add-on. It confirms R1-09: no group changed
its energy or CO2 per unit of value added clearly differently by 2018; exporters' ``27\% less energy'' was mostly size
(one clear row is left: their bought electricity, $-27$\%).}
\items{R1.S8, R1.S4, R1.S7, R1.Q1, R1.Q6, R1.G1, R1.G2, R1.G3}
\expart{Keep in mind}{Weights balance only 2013--14 averages. The intervals treat the weights as known (redraws that re-estimate
them change widths by $-8$\% to +12\%). Fictional firm links; association, not the effect of Paris.}
